\documentclass[twocolumn]{aastex631}

\usepackage{amsmath}
\usepackage{natbib}
\usepackage{graphicx}
\usepackage{multirow}

\shorttitle{Calibration Audit of Transient Classifiers}
\shortauthors{Sati}

\begin{document}

\title{Calibration Audit of Production Astronomical Transient Classifiers on ZTF Alert Streams}

\author{Pallavi Sati}
\affiliation{Independent Researcher}
\email{pallavisati23@gmail.com}

\begin{abstract}
Production transient classifiers deployed on the Zwicky Transient Facility (ZTF) alert stream are routinely used to prioritise spectroscopic follow-up, yet no prior work has evaluated calibration across multiple production brokers using standard metrics on real survey data. We present a systematic calibration audit of four classifiers: two production alert brokers (ALeRCE's Balanced Random Forest, and Fink's Random Forest and SuperNNova) and the standalone NEEDLE classification framework, using spectroscopically confirmed transients from the ZTF Bright Transient Survey as ground truth. We compute Expected Calibration Error (ECE) with equal-mass binning, classwise reliability diagrams, and Brier scores, then apply temperature scaling with five-fold cross-validation. We find that all four classifiers exhibit significant miscalibration. ALeRCE is systematically underconfident (ECE $= 0.271$); temperature scaling ($T = 0.357$) reduces ECE to $0.097$ (65\% improvement), increasing objects exceeding $p > 0.8$ from 20 to 425 at 91\% precision. Fink's Random Forest returns exactly zero scores for 94\% of objects, reflecting minimum-epoch data requirements rather than a classifier failure, rendering its outputs unsuitable as calibrated probabilities. SuperNNova's miscalibration (ECE $= 0.198$) is not addressable by temperature scaling. NEEDLE achieves good aggregate calibration at the object level (ECE $= 0.048$, $N = 278$ unique objects), but per-class analysis reveals severe class-imbalance-driven miscalibration: SLSN-I (superluminous supernovae type I, $n = 15$ true objects) are predicted with 73\% mean confidence but only 31\% accuracy (gap $= -0.423$), with three objects confidently misclassified as SNe at $> 0.998$ confidence. Global temperature scaling worsens calibration, and per-class scaling cannot resolve the structural miscalibration. These findings highlight the need for calibration-aware decision pipelines as the community prepares for LSST. Code and results are available at \url{https://github.com/pallavi-2000/transient-calibration-audit}.
\end{abstract}

\keywords{methods: statistical --- surveys --- supernovae: general --- techniques: photometric}

% ============================================================
\section{Introduction} \label{sec:intro}
% ============================================================

The Vera C. Rubin Observatory's Legacy Survey of Space and Time (LSST) will generate approximately 10 million alerts per night at full operations \citep{ivezic2019}, while global spectroscopic follow-up capacity remains fixed at roughly 100,000 spectra per year---a mismatch ratio of approximately 100:1. The first LSST alerts were issued in February 2026, with the full survey beginning later this year, making the calibration of broker classifiers an immediate operational concern rather than a future consideration. With the Zwicky Transient Facility \citep[ZTF;][]{bellm2019} having served as the primary precursor survey, the community has developed automated alert brokers to triage this data deluge and identify the most scientifically valuable transients for follow-up.

Several broker systems now operate in production on the ZTF alert stream, including ALeRCE \citep{forster2021}, Fink \citep{moller2021}, ANTARES \citep{narayan2018}, and Lasair \citep{smith2019}. These systems ingest alerts in real time, extract features from light curves and contextual information, and deploy machine learning classifiers to assign class probabilities to each transient. Astronomers routinely use these probability outputs to make follow-up decisions: a transient classified as ``90\% SN~Ia'' is far more likely to receive telescope time than one classified at ``60\% SN~Ia.'' Additional frameworks such as NEEDLE \citep{sheng2024} extend classification to rare transient types including superluminous supernovae (SLSNe) and tidal disruption events (TDEs).

This decision-making process relies on an implicit but critical assumption: that the classifier's stated probabilities are \emph{calibrated}---that is, among all objects assigned a probability $p$ of belonging to a given class, a fraction $p$ actually belongs to that class \citep{dawid1982}. Miscalibrated probabilities lead directly to misallocated resources: an overconfident classifier wastes telescope time on objects it is wrong about, while an underconfident one causes astronomers to overlook genuinely interesting transients. Despite the operational importance of calibration, no production transient classifier currently publishes calibration metrics such as the Expected Calibration Error \citep[ECE;][]{naeini2015}, reliability diagrams, or Brier score decompositions.

The machine learning calibration literature has established that modern neural networks are systematically overconfident \citep{guo2017}, a finding that has driven the development of post-hoc calibration methods including temperature scaling \citep{guo2017}, Platt scaling \citep{platt1999}, and Dirichlet calibration \citep{kull2019}. However, Random Forests---the architecture underlying several production broker classifiers---exhibit the opposite calibration behavior: ensemble averaging over bagged trees pushes predictions toward intermediate values, producing systematic underconfidence \citep{niculescumizil2005, bostrom2008}. Class imbalance corrections such as those used in Balanced Random Forests further distort probability estimates by inflating minority-class probabilities \citep{vandengoorbergh2022}. These well-characterized calibration properties of different architectures have direct relevance for astronomical classifiers but have not been evaluated in this domain.

Two recent works include limited calibration analysis. \citet{desoto2024} present per-class calibration curves for Superphot+ deployed on the ANTARES broker, noting overconfident SN~Ib/c and underconfident SN~Ia predictions but without computing formal calibration metrics. \citet{tung2025} evaluate calibration of ATCAT on simulated ELAsTiCC data, applying post-hoc corrections. However, no prior work has conducted a systematic calibration audit of production classifiers operating on real survey data using standard metrics.

In this paper, we present the first such audit. We evaluate four classifiers across three production systems: ALeRCE's Balanced Random Forest light curve classifier \citep[lc\_classifier v1.1.13;][]{sanchezsaez2021}, Fink's Random Forest and SuperNNova classifiers \citep{leoni2022, moller2020}, and the NEEDLE hybrid CNN+DNN framework \citep{sheng2024}. We restrict our analysis to classifiers that output explicit class probabilities; other brokers such as Lasair, ANTARES, and AMPEL employ crossmatching or feature-based approaches that do not produce probability estimates amenable to calibration analysis. Using spectroscopically confirmed transients from the ZTF Bright Transient Survey \citep[BTS;][]{fremling2020, perley2020} as ground truth, we compute ECE with equal-mass binning \citep{roelofs2022}, classwise reliability diagrams \citep{nixon2019}, and Brier scores, then apply temperature scaling with stratified five-fold cross-validation. We find that all four classifiers are significantly miscalibrated, each with a distinct failure mode, and that temperature scaling is effective for ALeRCE but inappropriate or counterproductive for the other systems.

The paper is organized as follows. Section~\ref{sec:background} reviews relevant background on calibration metrics and post-hoc methods. Section~\ref{sec:data} describes our data sources and sample construction. Section~\ref{sec:methods} presents our calibration analysis methodology. Section~\ref{sec:results} reports our findings for each classifier. Section~\ref{sec:discussion} interprets the results in the context of the calibration literature and discusses implications for LSST. Section~\ref{sec:conclusions} summarizes our conclusions.


% ============================================================
\section{Background} \label{sec:background}
% ============================================================

\subsection{Calibration and its Importance}

A probabilistic classifier is said to be \emph{calibrated} if, for all predicted probability values $p$, the empirical frequency of the positive class among predictions with confidence $p$ equals $p$ \citep{dawid1982}. Formally, for a classifier $f$ producing predictions $\hat{Y}$ with associated confidence $\hat{P}$:
\begin{equation}
P(\hat{Y} = Y \mid \hat{P} = p) = p \quad \forall p \in [0,1].
\end{equation}
Calibration is distinct from discrimination (the ability to separate classes). A classifier can achieve high accuracy while being poorly calibrated, and vice versa.

\subsection{Calibration Metrics}

The \emph{Expected Calibration Error} \citep[ECE;][]{naeini2015} is the most widely used scalar summary of calibration quality. It partitions predictions into $M$ bins based on confidence and computes the weighted average of the absolute difference between accuracy and confidence in each bin:
\begin{equation}
\text{ECE} = \sum_{m=1}^{M} \frac{|B_m|}{n} \left| \text{acc}(B_m) - \text{conf}(B_m) \right|,
\end{equation}
where $B_m$ is the set of samples in bin $m$, $\text{acc}(B_m)$ is the empirical accuracy, and $\text{conf}(B_m)$ is the mean confidence. We use $M = 15$ bins following \citet{guo2017}.

The choice of binning strategy significantly affects ECE estimates. Equal-width binning places bin boundaries at uniform intervals along the confidence axis, which can leave many bins empty or near-empty for classifiers with skewed confidence distributions. \citet{roelofs2022} demonstrated that equal-mass (quantile) binning, which places approximately the same number of samples in each bin, produces lower-bias ECE estimates and recommended against equal-width binning. We adopt equal-mass binning throughout.

Standard ECE considers only the top-predicted class and can mask miscalibration of individual classes. \citet{nixon2019} introduced the \emph{Static Calibration Error} (SCE) and \emph{Adaptive Calibration Error} (ACE), which evaluate calibration of all $K$ class probabilities independently. We report classwise ECE following this approach.

The \emph{Brier score} \citep{brier1950} provides a complementary proper scoring rule:
\begin{equation}
\text{BS} = \frac{1}{n} \sum_{i=1}^{n} \sum_{k=1}^{K} (p_{ik} - y_{ik})^2,
\end{equation}
where $p_{ik}$ is the predicted probability for class $k$ and $y_{ik}$ is the indicator for the true class.

\subsection{Temperature Scaling}

Temperature scaling \citep{guo2017} is a single-parameter post-hoc calibration method that rescales the logit vector $\mathbf{z}$ by a learned temperature $T > 0$:
\begin{equation}
\hat{q}_k = \frac{\exp(z_k / T)}{\sum_{j} \exp(z_j / T)}.
\end{equation}
When $T > 1$, the output distribution is softened (reducing overconfidence); when $T < 1$, it is sharpened (reducing underconfidence). The temperature is fitted by minimizing the negative log-likelihood (NLL) on a held-out calibration set, as NLL is a strictly proper scoring rule \citep{gneiting2007}.

Critically, temperature scaling applies a single global transformation and therefore cannot correct class-asymmetric miscalibration---where some classes are overconfident while others are underconfident. \citet{kull2019} demonstrated that temperature scaling can appear to produce well-calibrated confidence-level reliability diagrams while individual classes remain miscalibrated, and proposed Dirichlet calibration as a more flexible alternative.

\subsection{Calibration of Random Forests}

\citet{niculescumizil2005} showed that Random Forests produce systematically underconfident predictions due to the averaging of base learner probabilities: if the true probability is near 0 or 1, the ensemble average is biased toward 0.5 because only consensus among all trees can produce extreme probabilities. This effect is amplified by feature subsetting, which increases base learner variance. \citet{bostrom2008} confirmed this finding and proposed sigmoid calibration to correct it. \citet{loecher2024} demonstrated through simulation that Random Forest probability estimates ``are always underconfident'' across 192 scenarios.

Balanced Random Forests \citep{chen2004} introduce an additional calibration-distorting mechanism. By undersampling the majority class during training, they alter the effective class prior, which \citet{vandengoorbergh2022} showed produces ``strong and systematic overestimation of the probability for the minority class.'' The combined effect of standard RF underconfidence and minority-class inflation from balancing creates a complex but predictable calibration profile.

\subsection{Effects of Class Weighting on Calibration}

Many astronomical classifiers employ inverse-frequency class weighting or similar imbalance correction strategies to improve minority-class recall. \citet{vandengoorbergh2022} demonstrated that all common imbalance corrections---including random undersampling, oversampling, and SMOTE---harm probability calibration. \citet{fernando2021} confirmed that ``probability estimates of minority classes are unreliable'' when class weighting is applied. This is directly relevant to NEEDLE, which employs inverse-frequency weighting with extreme ratios ($\sim$80:1 for common supernovae versus TDEs).


% ============================================================
\section{Data} \label{sec:data}
% ============================================================

\subsection{Spectroscopic Ground Truth}

We use the ZTF Bright Transient Survey \citep[BTS;][]{fremling2020, perley2020} as our primary source of spectroscopic ground truth. BTS provides systematic spectroscopic classification of transients detected by ZTF, yielding 7,167 objects with spectroscopic type assignments. We map the $\sim$30 BTS spectroscopic subtypes to the ALeRCE taxonomy using a conservative mapping: SN~Ia and its subtypes (Ia-91T, Ia-91bg, Ia-02cx, Ia-CSM) map to SNIa; SN~Ib, Ic, Ic-BL, Ib/c, and Ibn map to SNIbc; SN~II, IIn, IIb, and IIP map to SNII; SLSN-I and SLSN-II map to SLSN; and TDE maps to TDE. Objects with types not present in any classifier taxonomy (novae, LRN, LBV) are excluded, yielding 7,116 mapped objects.

\subsection{Stratified Sampling}

Because rare classes (SLSN, TDE) constitute $<$2\% of the BTS catalog, we construct a stratified sample to ensure adequate representation for per-class calibration analysis. We cap the dominant class (SN~Ia) at 600 objects and include all available objects for rarer classes. This yields a sample of 1,436 objects with the distribution shown in Table~\ref{tab:sample}. Stratified sampling is appropriate for calibration analysis because ECE is a conditional measurement---it evaluates accuracy within confidence bins---and does not require the evaluation set to match the population class distribution \citep{guo2017}.

\begin{deluxetable}{lcc}
\tablecaption{Stratified Sample Composition \label{tab:sample}}
\tablehead{
\colhead{Class} & \colhead{Count} & \colhead{Fraction}
}
\startdata
SN~Ia & 600 & 41.8\% \\
SN~II & 341 & 23.7\% \\
SN~Ibc & 225 & 15.7\% \\
SLSN & 97 & 6.8\% \\
TDE & 39 & 2.7\% \\
Other & 134 & 9.3\% \\
\enddata
\end{deluxetable}

\subsection{Classifier Predictions}

\textbf{ALeRCE.} We query the ALeRCE API\footnote{\url{https://api.alerce.online}} for light curve classifier (lc\_classifier v1.1.13) probabilities. This Balanced Random Forest outputs a 15-class probability vector spanning transient and variable star categories \citep{sanchezsaez2021}. We successfully retrieve predictions for 1,149 of 1,436 sample objects. Since ALeRCE does not include a TDE class in its production taxonomy, we exclude 35 TDE objects, yielding 1,114 objects for analysis. We restrict to the four transient classes (SN~Ia, SN~Ibc, SN~II, SLSN) and renormalize the probability vectors to sum to unity over these classes. To verify that this renormalization does not introduce artifacts, we compared ECE on the renormalized 4-class vector (0.271) against the raw 15-class vector (0.283) and found the renormalized ECE is slightly lower. The median probability mass assigned to transient classes is 0.996, with only 1.4\% of objects (16/1,114) having $>$50\% probability in non-transient classes. Per-class inflation factors are negligible (1.02--1.06$\times$). Temperature scaling produces consistent results in both representations (post-calibration ECE of 0.097 vs.\ 0.092), confirming robustness to the renormalization choice. The absence of TDE from ALeRCE's production taxonomy is itself noteworthy: TDEs are a primary LSST science case, yet the most widely used ZTF transient classifier cannot identify them. \citet{pavezherrera2025} have recently added TDE support in a beta classifier, but it is not yet deployed in production.

\textbf{Fink.} We query the Fink API\footnote{\url{https://api.fink-portal.org}} for two binary classifiers: the Random Forest SN~Ia classifier (rf\_snia\_vs\_nonia) and SuperNNova \citep[snn\_snia\_vs\_nonia;][]{moller2020}. Both output a single score $\in [0,1]$ representing $P(\text{SN~Ia})$. We retrieve the most recent alert for each object (corresponding to the most complete light curve), obtaining predictions for 1,237 objects.

\textbf{NEEDLE.} Unlike ALeRCE and Fink, NEEDLE \citep{sheng2024} does not provide a public API. We perform local inference using the five pre-trained models from the \texttt{lasair\_th\_r} model family, applied to the publicly available dataset. Each model was trained with a different random held-out test set recorded in \texttt{testset\_obj.json}. We map test set ZTF IDs to dataset positions via the provided hash table, yielding 429 predictions across 278 unique objects in three classes (SN, SLSN-I, TDE). Of these, 43 objects appear in the held-out test sets of multiple models, contributing 151 duplicate predictions. Our primary analysis deduplicates to the object level by averaging model probabilities per object (278 unique objects); the 429-prediction model-instance set is retained as a sensitivity analysis (Section~\ref{sec:needle_sensitivity}). The 100\% inter-model class agreement (Section~\ref{sec:needle_sensitivity}) validates the averaging methodology.


% ============================================================
\section{Methods} \label{sec:methods}
% ============================================================

\subsection{Calibration Evaluation}

For each classifier, we compute:
\begin{enumerate}
    \item ECE with equal-mass binning (15 bins) and, for comparison, equal-width binning.
    \item Classwise ECE following \citet{nixon2019}, evaluating calibration of each class probability independently.
    \item Reliability diagrams showing empirical accuracy versus mean confidence per bin.
    \item Brier score as a proper scoring rule.
    \item Bootstrap confidence intervals (1,000 replicates, 95\% level) on ECE.
\end{enumerate}

For binary classifiers (Fink RF and SuperNNova), confidence is defined as $\max(p, 1-p)$ and predictions are thresholded at $p = 0.5$. We verified that ECE is robust to the choice of bin count for our smallest sample (NEEDLE, $N = 429$): ECE varies from 0.066 at $M = 10$ to 0.075 at $M = 20$, a range of 0.009.

\subsection{Temperature Scaling}

We fit the temperature parameter $T$ by minimizing negative log-likelihood using bounded scalar optimization ($T \in [0.01, 10.0]$). To avoid fitting and evaluating on the same data, we use stratified five-fold cross-validation: folds are constructed to preserve class ratios (essential when rare classes such as SLSN constitute only 7.6\% of the sample), $T$ is fitted on four folds, and ECE is evaluated on the held-out fold, with results averaged across folds.

For Random Forest classifiers, which do not natively produce logits, we convert vote-fraction probabilities to pseudo-logits via $z_k = \log p_k$ before applying temperature scaling. We note that this conversion assumes a log-linear relationship that may not hold precisely for RF outputs.

We apply three decision rules to determine whether temperature scaling should be recommended:
\begin{enumerate}
    \item \textbf{Bound detection:} If $T$ reaches the optimizer bound ($T \geq 9.5$ or $T \leq 0.02$), scaling is rejected as inappropriate.
    \item \textbf{Worsening detection:} If post-scaling ECE exceeds pre-scaling ECE, scaling is rejected.
    \item \textbf{Redundancy detection:} If pre-scaling ECE $< 0.1$, scaling is deemed unnecessary.
\end{enumerate}

For classifiers where global temperature scaling fails, we additionally evaluate per-class temperature scaling \citep{fang2020}, which fits an independent $T_k$ for each class and re-normalizes via softmax.


% ============================================================
\section{Results} \label{sec:results}
% ============================================================

Table~\ref{tab:summary} summarizes the calibration metrics for all four classifiers. We describe each in detail below.

\begin{deluxetable*}{lccccccc}
\tablecaption{Calibration Audit Summary \label{tab:summary}}
\tablehead{
\colhead{Classifier} & \colhead{Architecture} & \colhead{Classes} & \colhead{$N$} & \colhead{ECE} & \colhead{Brier} & \colhead{Post-hoc Method} & \colhead{ECE$_\mathrm{post}$}
}
\startdata
ALeRCE & BRF & 15 (4 transient) & 1114 & $0.271$ $[0.249, 0.296]$ & 0.486 & Temp.\ ($T=0.357$) & $0.097$ \\
Fink RF & RF & 2 & 1237 & $0.411$ $[0.382, 0.437]$ & 0.831 & N/A (degenerate) & --- \\
Fink SNN & RNN & 2 & 1237 & $0.198$ $[0.173, 0.226]$ & 0.513 & N/A (bound hit) & --- \\
\hline
NEEDLE (obj.-level)$^{\ddagger}$ & CNN+DNN & 3 & 278 & $0.048$ $[0.048, 0.111]$ & 0.265 & Global $T$ worsens & $0.169$ \\
\quad (model-instance) & & & 429 & $0.075$ $[0.062, 0.115]$ & 0.241 & (sensitivity; \S\ref{sec:needle_sensitivity}) & --- \\
\enddata
\tablecomments{ECE computed with equal-mass binning (15 bins). 95\% confidence intervals via bootstrap (1,000 replicates). Post-hoc ECE for ALeRCE is the cross-validated estimate (stratified 5-fold). Brier score: 0 is perfect; for reference, random guessing gives $\sim$0.5 for binary classifiers and $\sim$1.5 for 4-class. $^{\ddagger}$NEEDLE object-level (primary): 278 unique objects, deduplicated by averaging probabilities across models. Model-instance (sensitivity): all 429 predictions from 5 held-out test sets, treating each prediction independently. See \S\ref{sec:needle} and \S\ref{sec:robustness}.}
\end{deluxetable*}

\subsection{ALeRCE: Systematic Underconfidence} \label{sec:alerce}

The ALeRCE light curve classifier exhibits systematic underconfidence across all transient classes (Figure~\ref{fig:alerce_reliability}). The overall accuracy is 0.742, but the mean confidence is only 0.471---a gap of $+0.271$ indicating that the classifier is substantially more accurate than its confidence scores suggest. The reliability diagram shows all bins lying above the diagonal, with confidence values compressed into the range $[0.3, 0.75]$ and never exceeding 0.75.

This underconfidence is consistent with the known calibration behavior of Random Forests \citep{niculescumizil2005, bostrom2008}. The Balanced Random Forest architecture amplifies this effect: by undersampling the majority class during training, it alters the effective prior, and the averaging of base learner predictions across a 15-class taxonomy further pushes probabilities toward uniform values.

Per-class analysis (Figure~\ref{fig:alerce_perclass}) reveals that underconfidence is pervasive: SN~Ia (ECE $= 0.195$, gap $= +0.335$), SN~Ibc (ECE $= 0.154$, gap $= +0.201$), and SN~II (ECE $= 0.150$, gap $= +0.393$) all exhibit strong underconfidence. SLSN is a notable exception with a slight overconfidence (gap $= -0.055$), potentially reflecting the Balanced Random Forest's inflation of minority-class probabilities \citep{vandengoorbergh2022}.

Temperature scaling with $T = 0.357 \pm 0.012$ (stratified five-fold CV) reduces the cross-validated ECE from 0.276 to 0.097---a 65\% improvement (Figure~\ref{fig:alerce_tempscaling}). The sharpening effect ($T < 1$) is the correct direction for underconfident classifiers, spreading the compressed confidence range $[0.3, 0.75]$ out to $[0.35, 1.0]$. We note that $T = 0.357$ is more aggressive than values typically reported in the neural network calibration literature (where $T > 1$ is standard for overconfident models), reflecting the severity of RF underconfidence when distributed across 15 classes. The non-cross-validated ECE of 0.030 shown in Figure~\ref{fig:alerce_tempscaling} illustrates the visual effect of calibration but overestimates improvement due to fitting and evaluating on the same data; the cross-validated estimate of 0.097 is the appropriate value for reporting.

\begin{figure}
\plotone{fig1_alerce_reliability.pdf}
\caption{Reliability diagram for the ALeRCE light curve classifier showing systematic underconfidence: all bins lie above the perfect calibration diagonal. Confidence values are compressed into the range $[0.3, 0.75]$. ECE $= 0.271$ with equal-mass binning (15 bins). \label{fig:alerce_reliability}}
\end{figure}

\begin{figure*}
\plotone{fig2_alerce_perclass.pdf}
\caption{Per-class reliability diagrams for the ALeRCE classifier. SN~Ia, SN~Ibc, and SN~II show consistent underconfidence (bars above diagonal). SLSN shows slight overconfidence, potentially due to minority-class probability inflation from the Balanced Random Forest architecture. \label{fig:alerce_perclass}}
\end{figure*}

\begin{figure*}
\plotone{fig3_alerce_tempscaling.pdf}
\caption{ALeRCE reliability diagrams before (left) and after (right) temperature scaling with $T = 0.357$. The displayed ECE of 0.030 is computed on the full dataset; the cross-validated estimate is 0.097 (65\% improvement). The sharpening effect spreads the compressed confidence distribution to span the full $[0.35, 1.0]$ range. \label{fig:alerce_tempscaling}}
\end{figure*}


\subsection{Fink: Structural Limitations} \label{sec:fink}

\subsubsection{Random Forest}

Fink's Random Forest SN~Ia classifier produces degenerate probability outputs: 94\% of scores (1,161 of 1,237) are exactly 0.0 (Figure~\ref{fig:fink_rf}). The mean score is 0.011, and the formal ECE is 0.411. However, this ECE is misleading because the reliability diagram is dominated by a single bin at confidence $\approx 1.0$ (since $\max(0.0, 1 - 0.0) = 1.0$), containing nearly all objects.

This behavior is consistent with the classifier's design. \citet{leoni2022} describe a minimum data requirement of at least three observed epochs in each photometric filter before the classifier produces a score. Objects not meeting this threshold receive no classification, which the API returns as a score of 0.0. The high zero rate therefore likely reflects a mismatch between our sample (which includes objects at all evolutionary stages) and the classifier's operational constraints, rather than a fundamental architectural flaw. We classify this output distribution as unsuitable for probability calibration and recommend that Fink's RF scores be interpreted as detection flags rather than calibrated probabilities.

\subsubsection{SuperNNova}

Fink's SuperNNova classifier \citep{moller2020} implements a data-quality abstention gate, returning score $= 0$ for objects with insufficient photometric data. In our sample, 444 of 1,237 objects (35.9\%) trigger this gate, yielding 64.1\% effective prediction coverage (793 objects). The abstention rate is class-dependent: SN~Ia objects abstain at 19.0\% (100/527), while non-SN~Ia objects abstain at 48.5\% (344/710)---a 2.5$\times$ ratio consistent with sparser light curves for non-thermonuclear transients. This class asymmetry confirms that zeros represent deliberate abstentions rather than a data loading artefact. The first non-zero score is 0.0051, creating a hard gap from exactly zero and diagnostic of a threshold-based eligibility test (Figure~\ref{fig:fink_snn_hist}).

Among the 793 non-zero predictions, scores are actively discriminating (mean $= 0.564$, median $= 0.664$). The reliability diagram (Figure~\ref{fig:fink_snn}) shows bars consistently below the diagonal, indicating overconfidence (ECE $= 0.198$ [0.173, 0.226]).

Temperature scaling optimization drives $T$ to the upper bound ($T = 10.0$), indicating the method is inappropriate for this distribution. At $T = 10$ the optimizer attempts to flatten all predictions toward 0.5, destroying discriminative information. The combination of 35.9\% exact zeros (which temperature scaling cannot modify, as $\log(0) = -\infty$) and a continuous distribution for the remaining 64\% creates a bimodal structure incompatible with single-scalar calibration.

\begin{figure}
\plotone{fig_fink_random_forest.pdf}
\caption{Reliability diagram for Fink's Random Forest classifier. The degenerate output distribution (94\% of scores at exactly 0.0) produces a misleading reliability diagram concentrated at confidence $\approx 1.0$. ECE $= 0.411$. \label{fig:fink_rf}}
\end{figure}

\begin{figure}
\plotone{fig_fink_supernnova.pdf}
\caption{Reliability diagram for Fink's SuperNNova classifier (conditional on non-zero predictions, $n = 793$). Bars consistently below the diagonal indicate systematic overconfidence. ECE $= 0.198$. Temperature scaling hits the optimizer bound at $T = 10$. \label{fig:fink_snn}}
\end{figure}

\begin{figure}
\plotone{fig_fink_snn_zero_distribution.pdf}
\caption{Zero-score distribution for Fink SuperNNova. Left: full score distribution with 444 objects (35.9\%) at exactly zero and a hard gap of 0.005 to the first non-zero score, diagnostic of a data-quality abstention gate. Centre: abstention rate by true class---non-SN~Ia objects abstain at 2.5$\times$ the rate of SN~Ia objects (48.5\% vs.\ 19.0\%), confirming the zeros are class-dependent and not a loading artefact. Right: non-zero score distribution for SN~Ia and non-SN~Ia objects, showing discriminative power when predictions are made. \label{fig:fink_snn_hist}}
\end{figure}


\subsection{NEEDLE: Object-Level Calibration and Class Imbalance} \label{sec:needle}

\subsubsection{Aggregate Calibration: Object-Level Results}

Our primary NEEDLE analysis evaluates 278 unique ZTF objects, deduplicated by averaging probabilities from all models that classified each object (Section~\ref{sec:robustness}). At the object level, NEEDLE achieves aggregate ECE $= 0.048$ [95\% CI: 0.048, 0.111] and overall accuracy 80.2\% (Brier $= 0.265$). This aggregate ECE lies at the boundary of the ``well-calibrated'' threshold (ECE $< 0.05$), a markedly better result than the model-instance ECE of 0.075 that inflates rare-class contributions due to repeated appearances across held-out sets (Section~\ref{sec:needle_sensitivity}).

Figure~\ref{fig:needle_aggregate} shows the object-level reliability diagram. The aggregate bars track the diagonal reasonably well, consistent with the low ECE. However, as the per-class analysis below reveals, this aggregate view conceals severe, class-specific miscalibration driven by class imbalance.

\subsubsection{Per-Class Miscalibration and Class Imbalance} \label{sec:needle_perclass}

Table~\ref{tab:needle_perclass} and Figure~\ref{fig:needle_perclass} show the classwise calibration decomposition at the object level. The three classes exhibit fundamentally different calibration behavior:

\begin{deluxetable}{lcccccc}
\tablecaption{NEEDLE Per-Class Calibration (Object-Level, $N=278$) \label{tab:needle_perclass}}
\tablehead{
\colhead{Class} & \colhead{$n_\mathrm{true}$} & \colhead{$n_\mathrm{pred}$} & \colhead{Accuracy} & \colhead{Mean Conf.} & \colhead{Gap} & \colhead{Direction}
}
\startdata
SN     & 242 & 193 & 0.984 & 0.865 & $+0.119$ & Underconfident \\
SLSN-I &  15 &  39 & 0.308 & 0.730 & $-0.423$ & Overconfident  \\
TDE    &  21 &  46 & 0.457 & 0.752 & $-0.296$ & Overconfident  \\
\enddata
\tablecomments{Gap = Accuracy $-$ Mean Confidence (when class is predicted). Positive gap: classifier is more accurate than confident (underconfident). Negative gap: classifier is more confident than accurate (overconfident). $n_\mathrm{pred}$ exceeds $n_\mathrm{true}$ for SLSN-I and TDE because many non-rare objects are misclassified into those categories.}
\end{deluxetable}

\textbf{SN (majority class, $n = 242$).} SN predictions are underconfident: when NEEDLE predicts SN, it achieves 98.4\% accuracy but assigns only 86.5\% mean confidence (gap $= +0.119$). This is consistent with the systematic suppression of majority-class probabilities under inverse-frequency weighting, which upweights rare class gradients at the expense of the dominant class.

\textbf{SLSN-I (rare class, $n = 15$).} SLSN-I predictions are severely overconfident: when NEEDLE predicts SLSN-I, it achieves only 30.8\% accuracy but assigns 73.0\% mean confidence (gap $= -0.423$). Critically, 39 objects are predicted as SLSN-I despite only 15 being truly SLSN-I---meaning 24 of these predictions are false positives issued with high confidence. This is the dominant failure mode of the classifier.

\textbf{TDE (rare class, $n = 21$).} TDE predictions are similarly overconfident: 45.7\% accuracy against 75.2\% mean confidence (gap $= -0.296$). Of 46 TDE predictions, 25 are false positives. The pattern is analogous to SLSN-I, consistent with the same inverse-frequency weighting mechanism.

The asymmetry---SN underconfident while both rare classes are overconfident---is the canonical signature of inverse-frequency class weighting identified by \citet{vandengoorbergh2022}. The aggregate ECE of 0.048 conceals this because the majority-class underconfidence (SN, $n = 242$) partially cancels the minority-class overconfidence (SLSN-I + TDE, $n = 36$) in the weighted average.

\subsubsection{Silent Failures and Confidence-Accuracy Mismatch} \label{sec:needle_silent}

At the object level, five objects are classified with $\geq$90\% confidence but are incorrect (silent failures, 1.8\% of 278 objects). Table~\ref{tab:needle_silent} lists these objects.

\begin{deluxetable}{lccc}
\tablecaption{NEEDLE Object-Level Silent Failures ($\geq$90\% confidence, wrong) \label{tab:needle_silent}}
\tablehead{
\colhead{ZTF ID} & \colhead{Predicted} & \colhead{True} & \colhead{Confidence}
}
\startdata
ZTF20aagikvv & SN    & SLSN-I & 1.000 \\
ZTF21aakjkec & SN    & SLSN-I & 0.9997 \\
ZTF19abzoyeg & SN    & SLSN-I & 0.999 \\
ZTF21aagnwkk & TDE   & SN     & 0.901 \\
ZTF21aautijg & SLSN-I & SN    & 0.907 \\
\enddata
\tablecomments{Three of five silent failures are true SLSN-I objects misclassified as SN at near-certainty confidence. This pattern is consistent with the class-imbalance mechanism: rare objects are predicted as majority class with extreme confidence when they fall outside the learned minority-class decision region.}
\end{deluxetable}

Three of the five silent failures (ZTF20aagikvv, ZTF21aakjkec, ZTF19abzoyeg) are true SLSN-I objects classified as SN with confidence $> 0.998$. These objects fall outside NEEDLE's learned SLSN-I decision boundary despite their spectroscopic confirmation; the inverse-frequency weighting inflates the SLSN-I probabilities for objects \emph{within} the boundary while leaving ambiguous cases assigned entirely to the SN class with extreme confidence. For spectroscopic triage, a near-certainty SN prediction from NEEDLE carries an embedded risk: three independently spectroscopically confirmed SLSN-I objects were classified this way.

\subsubsection{Model-Instance Sensitivity: Robustness Across Held-Out Models} \label{sec:needle_sensitivity}

As a sensitivity analysis, we evaluate all 429 model-instance predictions, treating each model's prediction on its held-out test objects as independent (see Section~\ref{sec:robustness} for the deduplication methodology). The model-instance aggregate ECE is 0.075 [0.062, 0.115], with 17 silent failures (4.0\%).

The higher model-instance ECE (0.075 vs.\ 0.048 object-level) reflects an artifact of the evaluation design: rare-class objects (SLSN-I, TDE) appear in many more models' held-out test sets than SN objects (mean 5.0 models per rare object vs.\ 1.03 for SN), inflating the effective rare-class sample weight and shifting the aggregate ECE upward. The 100\% inter-model class agreement across all 43 multi-model objects---zero objects show disagreement in predicted class across models---validates that the averaging methodology used in the primary analysis is sound.

\subsubsection{Interpretation: Class Imbalance and Learned Biases} \label{sec:needle_interp}

The class-asymmetric miscalibration we observe is a direct consequence of NEEDLE's training procedure. \citet{sheng2024} employ inverse-frequency class weighting to address the extreme imbalance between common supernovae and rare transients, with approximately 80$\times$ more SN training examples than TDE examples. \citet{vandengoorbergh2022} demonstrated that such corrections systematically inflate minority-class probability estimates, which is precisely the pattern we observe for both SLSN-I (gap $= -0.423$) and TDE (gap $= -0.296$).

Global temperature scaling with $T = 1.552$ \emph{worsens} ECE from 0.073 to 0.169 (model-instance). This failure is expected: a single $T$ cannot simultaneously soften overconfident SLSN-I/TDE predictions ($T > 1$ needed) and sharpen underconfident SN predictions ($T < 1$ needed). \citet{kull2019} documented this fundamental limitation. Per-class temperature scaling yields $T_{\text{SN}} = 1.85$, $T_{\text{SLSN-I}} = 2.03$, and $T_{\text{TDE}} = 0.54$, confirming the opposing calibration directions, but worsens every individual classwise ECE---indicating structural miscalibration incompatible with post-hoc probability rescaling.

\begin{figure}
\plotone{fig_needle_dedup_reliability.pdf}
\caption{Aggregate reliability diagram for NEEDLE at the object level ($N = 278$ unique objects). ECE $= 0.048$ places NEEDLE at the boundary of the ``well-calibrated'' threshold, but this aggregate view conceals severe per-class miscalibration (Table~\ref{tab:needle_perclass} and Figure~\ref{fig:needle_perclass}). \label{fig:needle_aggregate}}
\end{figure}

\begin{figure*}
\plotone{fig_needle_perclass.pdf}
\caption{Per-class reliability diagrams for NEEDLE revealing class-asymmetric miscalibration hidden by the aggregate ECE. SN is underconfident (bars above diagonal); SLSN-I is severely overconfident (high-confidence bars below diagonal, especially near 1.0); TDE is overconfident. This asymmetry arises directly from inverse-frequency class weighting during training. \label{fig:needle_perclass}}
\end{figure*}

\begin{figure*}
\plotone{fig_needle_comparison_dedup_vs_full.pdf}
\caption{Comparison of NEEDLE calibration at model-instance level ($N = 429$) and object level ($N = 278$, primary). Left: aggregate ECE with 95\% bootstrap confidence intervals. Right: per-class ECE for SN, SLSN-I, and TDE. The object-level ECE is lower in aggregate because rare-class objects (which are more miscalibrated) appear in fewer models' test sets when counted once per object. \label{fig:needle_comparison}}
\end{figure*}

\subsection{Comparative Overview}

Figure~\ref{fig:summary} presents the ECE comparison across all four classifiers. At the object level, NEEDLE (ECE $= 0.048$) is the only classifier to reach the ``well-calibrated'' threshold (ECE $< 0.05$) in its aggregate form, but this masks severe per-class miscalibration for the rare transient classes. ALeRCE's temperature-scaled ECE of 0.097 falls below the ``acceptable'' threshold (ECE $< 0.10$). The remaining classifiers---Fink RF, Fink SNN, and NEEDLE at the class level---require either architectural changes or class-aware calibration approaches.

\begin{figure*}
\plotone{fig_summary_ece_comparison.pdf}
\caption{ECE comparison across all four production classifiers. NEEDLE (object-level, ECE $= 0.048$) is the only classifier to meet the ``well-calibrated'' threshold (dotted line), but this conceals class-level miscalibration for SLSN-I and TDE. Only ALeRCE, after temperature scaling, meets the ``acceptable'' threshold (ECE $< 0.10$, dashed line). Error bars show 95\% bootstrap confidence intervals. \label{fig:summary}}
\end{figure*}


% ============================================================
\section{Discussion} \label{sec:discussion}
% ============================================================

\subsection{Architecture Determines Calibration Failure Mode}

Our results demonstrate a clear relationship between classifier architecture and the nature of miscalibration. The ALeRCE Balanced Random Forest exhibits underconfidence consistent with the well-established calibration properties of ensemble methods \citep{niculescumizil2005, bostrom2008, loecher2024}. This is distinct from the overconfidence typically found in deep neural networks \citep{guo2017}, highlighting the importance of architecture-specific calibration analysis rather than assuming universal overconfidence.

NEEDLE's class-asymmetric miscalibration represents a qualitatively different failure mode that is invisible to standard aggregate ECE. This finding aligns with \citet{nixon2019}, who demonstrated that aggregate calibration metrics can conceal per-class miscalibration, and with \citet{kull2019}, who showed that confidence-level reliability diagrams can appear well-calibrated even when classwise diagrams reveal significant problems. We recommend that future calibration analyses of multi-class astronomical classifiers always include per-class evaluation.

\subsection{The Failure of Temperature Scaling as a Diagnostic}

The failure of global temperature scaling for NEEDLE and Fink is itself informative. When temperature scaling worsens calibration, it indicates that the miscalibration is structural---arising from the model's architecture or training procedure---rather than a simple uniform shift in confidence. For NEEDLE, the opposing per-class temperatures ($T_{\text{SN}} = 1.85$ versus $T_{\text{TDE}} = 0.54$) directly trace the effect of inverse-frequency class weighting, confirming the mechanism identified by \citet{vandengoorbergh2022}. For Fink SuperNNova, the bound-hitting behavior ($T = 10$) reveals that the bimodal score distribution (36\% exact zeros plus a continuous component) is fundamentally incompatible with temperature scaling's assumption of a smooth logit space.

These diagnostic patterns suggest a practical recommendation: apply temperature scaling not only as a correction method but as a probe of miscalibration structure. If $T \approx 1$, the classifier is approximately calibrated. If $T$ converges to a finite value away from 1 and ECE improves, the miscalibration is uniform and fixable. If $T$ hits a bound or worsens ECE, deeper architectural or training-level interventions are needed.

\subsection{Implications for Spectroscopic Follow-up}

Uncalibrated probabilities have direct, quantifiable consequences for follow-up allocation. Table~\ref{tab:decision} shows the effect of calibration on ALeRCE follow-up decisions at various probability thresholds.

\begin{deluxetable}{ccccc}
\tablecaption{Follow-up Decisions Before and After Calibration \label{tab:decision}}
\tablehead{
\colhead{Threshold} & \multicolumn{2}{c}{Raw (uncalibrated)} & \multicolumn{2}{c}{Calibrated ($T = 0.357$)} \\
\colhead{$p >$} & \colhead{$N$ pass} & \colhead{Precision} & \colhead{$N$ pass} & \colhead{Precision}
}
\startdata
0.5 & 392 & 0.898 & 972 & 0.784 \\
0.6 & 134 & 0.940 & 812 & 0.828 \\
0.7 & 40 & 1.000 & 627 & 0.871 \\
0.8 & 20 & 1.000 & 425 & 0.911 \\
0.9 & 1 & 1.000 & 194 & 0.948 \\
\enddata
\tablecomments{$N$ pass: number of objects exceeding the confidence threshold. Precision: fraction of those objects correctly classified. Based on 1,114 ALeRCE predictions with stratified 5-fold cross-validated temperature scaling.}
\end{deluxetable}

The operational impact is stark. At a $p > 0.8$ follow-up threshold---a reasonable requirement for committing telescope time---only 20 objects pass using raw probabilities, because ALeRCE's confidence never exceeds $\sim$0.75. After calibration, 425 objects pass at 91.1\% precision: a 21$\times$ increase in usable classifications.

For NEEDLE, the operational concern is SLSN-I overconfidence. When the classifier predicts SLSN-I, the mean confidence is 73.0\% but the true accuracy is only 30.8\% (per-class gap $= -0.423$), exemplifying how inverse-frequency class weighting shifts learned decision boundaries without adjusting confidence calibration. For spectroscopic triage, this means: an SLSN-I prediction carries only $\sim$31\% prior probability of being correct, far below the stated confidence. Astronomers relying on this classifier for rare-object discovery must understand this behavior. Three silently failing SLSN-I objects (ZTF20aagikvv, ZTF21aakjkec, ZTF19abzoyeg) are classified as SN at near-certainty confidence ($> 0.998$), meaning per-class reliability diagrams---not just aggregate ECE---are essential for operational trust.

\subsection{Implications for LSST}

With LSST's first alerts already issued and the full survey beginning in 2026, the calibration problems we identify in ZTF-era classifiers demand immediate attention. The greater diversity of transient types, the increased class imbalance (rarer classes become proportionally rarer in a larger survey), and the higher volume of alerts all exacerbate miscalibration risk. Our findings suggest three concrete recommendations:

First, broker teams should publish calibration metrics alongside classification performance. ECE, reliability diagrams, and Brier scores should be standard evaluation outputs, not afterthoughts.

Second, post-hoc calibration should be architecture-aware. Temperature scaling is appropriate for Random Forest classifiers but can be counterproductive for class-weighted neural networks. Class-aware methods such as Dirichlet calibration \citep{kull2019} or focal calibration \citep{mukhoti2020} may be more appropriate for classifiers with extreme class imbalance.

Third, calibration should be evaluated on real survey data, not only on simulated datasets. Our finding that Fink's RF produces 94\% zero scores when evaluated on BTS objects---behavior that would not appear in simulated data designed to match the classifier's training distribution---illustrates the gap between simulated and operational calibration.

\subsection{Robustness Checks} \label{sec:robustness}

We performed five robustness checks to verify that primary findings are not artefacts of data processing or methodological choices.

\textbf{ALeRCE class restriction.} Restricting ALeRCE's 15-class output to the 4-class transient subset does not drive the observed miscalibration: the 4-class renormalized ECE (0.271) matches the native 15-class ECE (0.283) within 4\%, with negligible per-class probability inflation (1.02--1.06$\times$). The 15-class analysis serves as a harder baseline (more classes, lower accuracy) confirming the 4-class restriction is conservative.

\textbf{NEEDLE object-level deduplication.} The NEEDLE evaluation set contains 429 model-instance predictions across 278 unique objects, with 43 objects appearing in two or more held-out test sets. Deduplicating to object level via mean-of-model-probabilities yields ECE $= 0.048$ (model-instance: 0.075), a 36\% reduction. The 100\% inter-model class agreement across all 43 multi-model objects (zero class disagreements) validates the averaging methodology: model variance in predicted class is zero, and confidence uncertainty dominates. Per-class miscalibration patterns persist in both analyses. Object-level results are primary; model-instance results are reported in \S\ref{sec:needle_sensitivity}.

\textbf{ALeRCE operational gain: held-out cross-validation.} The paper's 21$\times$ operational gain (20 $\to$ 425 objects at $p > 0.80$) was originally computed on the full dataset with temperature $T = 0.357$. To address potential data leakage, we recomputed with stratified 5-fold cross-validation: $T$ fitted on four folds, gain measured on the held-out fold. The pooled held-out result (20 $\to$ 418 objects, gain $= 20.9\times$, precision $= 0.907$) is essentially identical to the full-dataset claim, because $T$ is a single scalar that converges to $0.357 \pm 0.012$ regardless of which fold trains it. Per-fold gain shows high variance ($30\times \pm 23\times$) due to tiny pre-scaling counts per fold (1--6 objects), but the pooled estimate is the appropriate metric (Figure~\ref{fig:alerce_cv}).

\textbf{Calibration method comparison (Fink RF).} We compared temperature scaling, Platt scaling, and isotonic regression on Fink RF's 76 non-zero predictions via 5-fold cross-validation. Isotonic regression achieves the lowest point estimate (ECE $= 0.123 \pm 0.058$) versus temperature scaling (ECE $= 0.175 \pm 0.055$), but the high variance at $N = 76$ ($\approx$12 test objects per fold) makes the difference inconclusive. Temperature scaling is the more conservative choice at this sample size; isotonic regression represents a potential improvement if Fink RF coverage expands (Figure~\ref{fig:fink_methods}).

\textbf{Bin-count sensitivity.} ECE was computed at 5, 10, 15, and 20 equal-mass bins for all classifiers (Figure~\ref{fig:bins}). ALeRCE ($N = 1114$) shows zero relative variation (0.0\%), Fink SNN ($N = 793$) 6.3\%, and Fink RF ($N = 76$) 9.2\%, all below the ``acceptable'' threshold of 10\%. NEEDLE ($N = 278$, 3 classes) shows 56.5\% relative variation due to small sample size and rare-class imbalance: with 15 bins, only $\approx$18 objects occupy each bin, and the 15 SLSN-I and 21 TDE true objects span fractional bins. The wide CI $[0.048, 0.111]$ already captures this bin-count uncertainty; the 15-bin point estimate of 0.048 is neither the minimum nor maximum across bin counts (0.036 at 5--10 bins; 0.061 at 20 bins). For large-sample classifiers, 15 bins is well within the stable region.

\subsection{Limitations}

Several limitations should be noted. Our ground truth is limited to spectroscopically confirmed objects from BTS, which is itself magnitude-limited and subject to selection effects. The stratified sampling may not capture the full range of classifier behavior across the alert stream. For NEEDLE, our evaluation uses predictions from five models on their respective held-out test sets, which may not represent the classifier's performance on the broader alert stream.

The small sample sizes for rare classes at the object level ($n = 15$ SLSN-I, $n = 21$ TDE for NEEDLE) limit the statistical power of per-class calibration analysis. We emphasise that per-class reliability diagrams for these rare classes should be interpreted as indicating the \emph{direction} of miscalibration (overconfident vs.\ underconfident) rather than as precise ECE estimates. The wide confidence intervals on the aggregate NEEDLE ECE [0.048, 0.111] reflect this fundamental sample-size constraint. More flexible calibration methods such as isotonic regression or Dirichlet calibration require larger samples to fit reliably \citep{vaicenavicius2019}, which is why we do not attempt them here.

Our analysis uses the most recent alert for each Fink object, corresponding to the most complete light curve. In operational use, astronomers require classification decisions early in a transient's evolution, when light curves are incomplete and classifier confidence is lower. How calibration varies as a function of light curve completeness is a natural and important extension of this work, particularly for LSST where early classification will be critical.

The pseudo-logit transformation $z_k = \log p_k$ used to apply temperature scaling to Random Forest outputs lacks formal theoretical justification---temperature scaling was designed for softmax-based neural network outputs, not ensemble vote fractions. While the consistent cross-validated improvement (ECE reduction from 0.276 to 0.097 with $T$ variance of $\pm 0.012$) suggests the transformation is empirically effective, future work could compare alternative transformations such as Platt scaling applied directly to vote fractions.

Finally, even equal-mass ECE is a biased estimator of true calibration error \citep{roelofs2022}. We report bootstrap confidence intervals throughout to quantify estimation uncertainty. Bin-count sensitivity analysis (\S\ref{sec:robustness}) shows that ECE is stable for all classifiers at large sample sizes; for NEEDLE ($N = 278$) the wide CI $[0.048, 0.111]$ already subsumes bin-count uncertainty.


% ============================================================
\section{Conclusions} \label{sec:conclusions}
% ============================================================

We have presented the first systematic calibration audit of production machine learning classifiers deployed on ZTF transient alert streams. Our key findings are:

\begin{enumerate}
    \item No production transient classifier currently publishes calibration metrics, yet astronomers routinely use their probability outputs for follow-up decisions.

    \item ALeRCE's Balanced Random Forest is systematically underconfident (ECE $= 0.271$), consistent with known RF calibration behavior. Temperature scaling with $T = 0.357$ reduces ECE to 0.097 (65\% improvement), increasing the number of objects passing a $p > 0.8$ follow-up threshold from 20 to 425 at 91\% precision.

    \item Fink's Random Forest produces degenerate outputs (94\% zeros), likely reflecting minimum-epoch requirements rather than a fundamental flaw, but rendering its scores unsuitable as calibrated probabilities. Fink's SuperNNova is overconfident (ECE $= 0.198$) with miscalibration that temperature scaling cannot address.

    \item NEEDLE achieves aggregate ECE $= 0.048$ at the object level ($N = 278$ unique objects), meeting the well-calibrated threshold---but per-class analysis reveals severe class-imbalance-driven miscalibration. SLSN-I predictions are issued with 73\% mean confidence but only 31\% accuracy (gap $= -0.423$), and three true SLSN-I objects are classified as SN with near-certainty confidence ($> 0.998$). TDE predictions show similar overconfidence (gap $= -0.296$). Global temperature scaling worsens calibration, and per-class scaling worsens every individual class---indicating structural miscalibration that post-hoc methods cannot resolve.

    \item The failure mode of temperature scaling is itself diagnostic: worsening ECE indicates structural miscalibration requiring architectural intervention, not post-hoc correction.
\end{enumerate}

As time-domain astronomy enters the LSST era, calibrated probabilities are essential for optimal spectroscopic resource allocation. We recommend that broker teams adopt calibration evaluation as a standard component of classifier assessment, using proper scoring rules and per-class reliability analysis.


% ============================================================
\section*{Data and Code Availability}

All analysis code, data acquisition scripts, and results are publicly available at \url{https://github.com/pallavi-2000/transient-calibration-audit}. The repository includes scripts to reproduce all figures and tables from the raw data. Spectroscopic ground truth is drawn from the publicly accessible ZTF Bright Transient Survey \citep{fremling2020, perley2020}. ALeRCE and Fink predictions were obtained via their respective public APIs. NEEDLE predictions were generated using publicly available pre-trained models and datasets.


% ============================================================
\begin{acknowledgments}
We thank Matt Nicholl for encouragement and for suggesting connections with the NEEDLE team, and Xinyue Sheng for making the NEEDLE models and dataset publicly available. This work made use of the ALeRCE broker \citep{forster2021}, the Fink broker \citep{moller2021}, and data from the ZTF Bright Transient Survey \citep{fremling2020, perley2020}. This research has made use of NASA's Astrophysics Data System.
\end{acknowledgments}


% ============================================================
\software{
\texttt{NumPy} \citep{harris2020},
\texttt{SciPy} \citep{virtanen2020},
\texttt{pandas} \citep{mckinney2010},
\texttt{Matplotlib} \citep{hunter2007},
\texttt{TensorFlow} \citep{abadi2016}
}


% ============================================================
\bibliography{references}
\bibliographystyle{aasjournal}

% ============================================================
\appendix
% ============================================================

\section{Cross-Validated Operational Gain (ALeRCE)} \label{app:cv}

\begin{figure*}
\plotone{fig_alerce_operational_gain_cv.pdf}
\caption{Cross-validated operational gain for ALeRCE temperature scaling (5-fold stratified CV). \emph{Left}: temperature $T$ per fold; mean $T = 0.357 \pm 0.012$ is virtually identical to the full-dataset estimate, confirming that the scalar calibration parameter is insensitive to training-set composition and that no data leakage inflates the paper's operational gain claim. \emph{Centre}: per-fold and pooled gain factor at $p > 0.80$; the pooled held-out result ($20 \to 418$ objects, gain $= 20.9\times$) matches the paper claim of 21$\times$. High per-fold variance (15--76$\times$) arises because the pre-scaling count per fold is 1--6 objects. \emph{Right}: per-fold ECE before (0.276 $\pm$ 0.033) and after (0.097 $\pm$ 0.016) temperature scaling, consistent with the full-dataset improvement. \label{fig:alerce_cv}}
\end{figure*}


\section{Calibration Method Comparison for Fink RF} \label{app:methods}

\begin{figure*}
\plotone{fig_fink_rf_calibration_methods.pdf}
\caption{Comparison of three post-hoc calibration methods for Fink RF non-zero predictions ($N = 76$, 5-fold CV). Temperature scaling (ECE $= 0.175 \pm 0.055$), Platt scaling (ECE $= 0.198 \pm 0.096$), and isotonic regression (ECE $= 0.123 \pm 0.058$) all reduce ECE from the raw value of 0.324. Isotonic regression achieves the lowest point estimate but with the highest variance; the overlapping uncertainty ranges and risk of overfitting on $\approx$12 test objects per fold make temperature scaling the more conservative recommendation at $N = 76$. \label{fig:fink_methods}}
\end{figure*}


\section{Bin-Count Sensitivity Analysis} \label{app:bins}

\begin{figure*}
\plotone{fig_bin_sensitivity_analysis.pdf}
\caption{ECE as a function of bin count (5, 10, 15, 20 equal-mass bins) for all four classifiers. ALeRCE ($N = 1114$): 0.0\% relative variation (excellent). Fink RF ($N = 76$): 9.2\% (acceptable). Fink SNN ($N = 793$): 6.3\% (acceptable). NEEDLE ($N = 278$): 56.5\% (poor, driven by small sample and rare-class imbalance). The dashed vertical line marks the paper's choice of 15 bins; it is neither a minimum nor maximum for any classifier, and the choice is well within the stable region for ALeRCE and Fink SNN. For NEEDLE, the wide bootstrap CI $[0.048, 0.111]$ already subsumes the bin-count uncertainty. \label{fig:bins}}
\end{figure*}

\end{document}
